\chapter{Continuous and Discrete Signals}

In the previous chapter we agreed that a signal is just a rule that assigns a
value to a moment of time -- a function. Such functions come in two kinds, and
almost everything else in this book works with one or the
other. On one side: \emph{continuous-time} signals, the smooth
voltages and pressures and accelerations of the physical world, defined at every
instant. On the other: \emph{discrete-time} signals, the lists of numbers that a
microcontroller actually stores and processes, defined only at a clock's tick.

The gesture controller we are building lives on \emph{both} sides at once. Your
wrist moves through real, continuous space. An accelerometer feels that motion as
a continuous voltage. But the chip that decides ``that was a flick, send a MIDI
note'' only ever sees a stream of integers arriving 200 times a second. Much of
this book is about reasoning clearly across that boundary -- knowing what
sampling preserves and what it loses. This chapter builds the two worlds
side by side, gives you the small zoo of basic signals you will meet again in
every later chapter, and introduces the discrete impulse, the single-sample spike
from which we will build convolution and every digital filter in the book.

\chapterline

\section{Two Kinds of Time}
\label{ch02:two-kinds}

Picture an audio recording. The microphone's diaphragm is pushed back and forth
by air pressure; at every instant across the whole timeline -- at
\emph{every} real-numbered instant -- there is a definite pressure. Call it
$x(t)$. The independent variable $t$ ranges over the whole real line $\R$ (or some
interval of it). There is a value at $t = 1$ second, at $t = 1.5$, at
$t = 1.500001$, at $t = \pi/7$. Between any two instants there are infinitely many
more. This is a \textbf{continuous-time} signal.

Now picture the same sound after it has gone through the analog-to-digital
converter in your laptop. The converter looks at the pressure, writes down a
number, waits a fixed interval, looks again, writes down another number, and so
on -- 44{,}100 times per second for CD audio. What you are left with is a
\emph{sequence}: $x[0], x[1], x[2], \dots$ There is no such thing as
``$x[1.5]$''. The signal is only defined at integer indices. This is a
\textbf{discrete-time} signal.

\begin{focusbox}
The type of bracket matters. In this book, and in most of the field, we write
parentheses for continuous time, $x(t)$ with $t \in \R$, and square brackets for
discrete time, $x[n]$ with $n \in \Z$. When you see $x[n]$, the variable $n$ is an
\emph{integer index}: it counts samples, not seconds. The conversion between the
index $n$ and a real time is $t = n T_s$, where $T_s$ is the time between samples.
\end{focusbox}

\begin{definition}[Continuous- and discrete-time signals]
\label{ch02:defn-ct-dt}
A \emph{continuous-time} signal is a function $x : \R \to \R$ (or $\R \to \C$),
written $x(t)$, defined for every real $t$ in its domain.
A \emph{discrete-time} signal is a function $x : \Z \to \R$ (or $\Z \to \C$),
written $x[n]$, defined only for integer $n$. We call $x[n]$ a \emph{sequence}.
\end{definition}

A discrete-time signal is best drawn not as a curve but as a \emph{stem plot}: a
dot at each integer $n$, with a thin stalk connecting it to the axis, to remind
you that there is nothing in between. Figure~\ref{ch02:fig-ct-vs-dt} shows the
same underlying sinusoid in both worlds.

\begin{figure}[h]
\centering
\begin{tikzpicture}
\begin{axis}[
title style={font=\small},
  width=0.92\textwidth, height=5.2cm,
  xlabel={$t$ (seconds)}, ylabel={amplitude},
  xmin=-0.05, xmax=1.05, ymin=-1.4, ymax=1.4,
  axis lines=left, samples=160, domain=0:1,
  title={Continuous $x(t)=\cos(2\pi\cdot 3\,t)$ and its samples $x[n]=x(nT_s)$, $T_s=1/12$},
  title style={font=\small},
]
\addplot[sigblue, very thick] {cos(deg(2*pi*3*x))};
\addplot[specred, only marks, mark=*, mark size=1.6pt,
  samples at={0,0.0833,0.1667,0.25,0.3333,0.4167,0.5,0.5833,0.6667,0.75,0.8333,0.9167,1.0}]
  {cos(deg(2*pi*3*x))};
\foreach \k in {0,0.0833,0.1667,0.25,0.3333,0.4167,0.5,0.5833,0.6667,0.75,0.8333,0.9167,1.0} {
  \addplot[specred, thin] coordinates {(\k,0) (\k,{cos(deg(2*pi*3*\k))})};
}
\end{axis}
\end{tikzpicture}
\caption{The continuous signal (blue curve) exists at every instant. The discrete
signal (red stems) is the list of values read off at the sampling instants
$t = nT_s$. The discrete signal knows nothing about the curve between its dots.}
\label{ch02:fig-ct-vs-dt}
\end{figure}

\begin{intuitionbox}
A continuous signal is an unbroken profile; a discrete signal is a \emph{list of
elevations measured along a hike}. From the list alone you cannot say what the
ground did between your steps -- maybe it was smooth, maybe there was a hole you
stepped over. Sampling throws away the in-between, and a huge fraction of signal
processing is about knowing when that is safe (the sampling theorem) and when it
bites you (aliasing).
\end{intuitionbox}

\subsection*{Why we even bother with the continuous world}

You might object: computers are digital, so why not study only sequences? Because
the physics is continuous. The wrist moves smoothly; light, sound, and EMG
voltages are continuous; the laws that govern them (Newton, Maxwell) are
differential equations in continuous time. If you only ever think in samples you
will be forever surprised by phenomena -- aliasing, the need for anti-alias
filters, the meaning of ``frequency'' above half the sample rate -- that only make
sense when you remember the continuous signal that the samples came from. We keep
both views because many practical bugs arise at the boundary between them.

\begin{connectionbox}
This fork mirrors a fork in the companion calculus book: the integral
$\int x(t)\,\dd t$ for continuous accumulation versus the sum $\sum_n x[n]$ for
discrete accumulation. Almost every continuous-time formula in this book has a
discrete twin in which $\int \to \sum$ and $\dd t \to 1$ (or $T_s$). Energy,
convolution, the Fourier transform -- all come in these two flavors, and seeing
them as continuous and discrete versions of the same operation will save you a
great deal of memorization.
\end{connectionbox}

\section{Sampling, as a First Idea}
\label{ch02:sampling}

We will devote a whole later chapter to sampling done right. But the basic move
is so central that we introduce it now, in its simplest form, so the rest of this
chapter can speak both languages.

\textbf{Ideal sampling} means: pick a fixed spacing $T_s$ (the \emph{sampling
period}, in seconds), and define the sequence
\[
  x[n] \;=\; x(nT_s), \qquad n \in \Z .
\]
The reciprocal $f_s = 1/T_s$ is the \emph{sampling rate} or \emph{sampling
frequency}, measured in samples per second, or hertz (Hz). If $f_s = 200$~Hz then
$T_s = 5$~ms, and $x[n]$ is the value of the underlying signal $5n$ milliseconds
after we started.

\begin{example}[From seconds to indices]
\label{ch02:ex-index}
An accelerometer is sampled at $f_s = 200$~Hz. You want the sample nearest to
$t = 0.31$~s. Then $n = t f_s = 0.31 \times 200 = 62$. Conversely, sample
$n = 150$ corresponds to $t = n/f_s = 150/200 = 0.75$~s. Get comfortable sliding
between $t$ and $n$; you will do it constantly.
\end{example}

\begin{trapbox}
Sampling does \emph{not} mean ``rounding the values.'' Sampling discretizes
\emph{time}; rounding the values to a finite set of levels is a separate step
called \emph{quantization} (Chapter~19). A signal can be discrete in time but
have perfectly real-valued (infinitely fine) amplitudes -- that is exactly the
$x[n] \in \R$ of Definition~\ref{ch02:defn-ct-dt}. Real ADCs do both at once, but
they are different operations and confusing them causes real bugs.
\end{trapbox}

\begin{goalbox}
The wrist tracker's MPU-6050 IMU is, in effect, an array of tiny springs whose
deflection becomes a continuous voltage. The chip's built-in ADC samples that
voltage, typically at something between 100~Hz and 1~kHz, and hands your
microcontroller a sequence $a_x[n], a_y[n], a_z[n]$ of acceleration samples. Every
gesture-detection idea later in the book operates on these sequences. Choosing
$f_s$ is your first real design decision: too low and a fast flick of the wrist
slips between samples (aliasing, lost gestures); too high and you waste power and
flood the CPU. Human hand motion has very little energy above about 20~Hz, so an
$f_s$ of a few hundred hertz is plenty -- but we will only be able to
\emph{justify} that number once we have the sampling theorem.
\end{goalbox}

\section{The Small Zoo of Basic Signals}
\label{ch02:zoo}

A handful of elementary signals appear over and over. They are the ``primary
colors'' of signal processing: more complicated signals are built from them, and
the behavior of a whole system can often be summarized by how it responds to one
of them. We introduce each in both continuous and discrete form.

\subsection{The unit impulse}
\label{ch02:impulse}

In discrete time the impulse is the simplest signal imaginable: it is $1$ at the
origin and $0$ everywhere else.
\[
  \delta[n] \;=\;
  \begin{cases}
    1, & n = 0,\\
    0, & n \neq 0.
  \end{cases}
\]
This is the \emph{Kronecker delta} and there is nothing subtle about it -- a single
spike, as in the left panel of Figure~\ref{ch02:fig-basic}. Despite its
simplicity, by the end of this chapter (and decisively in Chapter~13)
it will turn out that \emph{every} discrete signal is a sum of shifted,
scaled impulses, and that a linear system is completely described by its response
to this one spike.

The continuous-time impulse $\delta(t)$ -- the \emph{Dirac delta} -- is a stranger
beast. We would like a signal that is ``zero everywhere except at $t=0$, yet has
total area $1$.'' No ordinary function does this: a single point has zero width, so
zero area. The Dirac delta is defined not by its values but by what it does inside
an integral, through the \emph{sifting} property
\[
  \int_{-\infty}^{\infty} x(t)\,\delta(t)\,\dd t \;=\; x(0),
\]
and more generally $\int x(t)\,\delta(t - t_0)\,\dd t = x(t_0)$. Think of
$\delta(t)$ as the limit of a tall narrow rectangle of width $\eps$ and height
$1/\eps$ as $\eps \to 0$: the area stays $1$ while the spike grows infinitely thin
and tall. We draw it as an upward arrow of ``height 1'' (really, of \emph{area} 1).

\begin{formalbox}
Strictly, $\delta(t)$ is not a function but a \emph{distribution} (generalized
function): a rule that eats a smooth test function $x(t)$ and returns the number
$x(0)$. Engineers manipulate it as if it were a function and almost never get into
trouble, because every legitimate use ends with $\delta$ inside an integral, where
the sifting property tells you the answer. The discrete $\delta[n]$ needs none of
this care -- it is a plain, finite sequence that is $1$ at zero and $0$ elsewhere.
\end{formalbox}

\subsection{The unit step}
\label{ch02:step}

The unit step ``turns on'' at the origin:
\[
  u[n] =
  \begin{cases} 1, & n \ge 0 \\ 0, & n < 0 \end{cases}
  \qquad\qquad
  u(t) =
  \begin{cases} 1, & t \ge 0 \\ 0, & t < 0. \end{cases}
\]
The step models switching something on: applying a constant force at $t=0$, closing
a circuit, the onset of a sustained note. It also lets you \emph{window} a signal:
$x[n]\,u[n]$ is $x[n]$ for $n \ge 0$ and zero before, which is how you describe a
signal that ``starts'' at a definite time.

Step and impulse are intimately related. In discrete time, the step is the running
sum of impulses, and the impulse is the first difference of the step:
\[
  u[n] = \sum_{k=-\infty}^{n} \delta[k], \qquad
  \delta[n] = u[n] - u[n-1].
\]
In continuous time the same pair appears with calculus instead of sums:
$u(t) = \int_{-\infty}^{t}\delta(\tau)\,\dd\tau$ and, formally,
$\delta(t) = \dfrac{\dd}{\dd t}u(t)$ -- the derivative of a jump is an impulse.

\begin{intuitionbox}
Summation/integration and differencing/differentiation are inverse operations, so
$\delta$ and $u$ form a derivative/antiderivative pair. This is your first glimpse
of a recurring theme: the continuous and discrete worlds run the same play with
$\int \leftrightarrow \sum$ and $\frac{\dd}{\dd t} \leftrightarrow$ (first
difference).
\end{intuitionbox}

\subsection{The unit ramp}

The ramp is the running sum (or integral) of the step: it rises linearly once it
switches on.
\[
  r[n] = n\,u[n], \qquad\qquad r(t) = t\,u(t).
\]
Concretely, $r[n] = \dots,0,0,0,1,2,3,\dots$ starting its climb at $n=0$. The trio
impulse $\to$ step $\to$ ramp is each the running accumulation of the one before
it; differencing walks back the other way. The ramp shows up whenever a quantity
grows at a constant rate after some onset -- a position under constant velocity, or
a phase that increases linearly (which we will soon see \emph{is} a
constant-frequency sinusoid).

\subsection{The real exponential}
\label{ch02:exp}

\[
  x(t) = e^{a t} \quad (a \in \R), \qquad\qquad x[n] = \alpha^{\,n} \quad (\alpha \in \R).
\]
For the continuous form, $a < 0$ gives exponential \emph{decay} (the discharge of
a capacitor, the die-away of a struck string), $a > 0$ gives growth, and $a = 0$
gives a constant. For the discrete form $\alpha^n$: $\abs{\alpha} < 1$ decays,
$\abs{\alpha} > 1$ grows, and a \emph{negative} $\alpha$ makes the sign alternate
each sample -- a flavor of behavior with no continuous-time analog among real
exponentials, and a hint that the fastest wiggle a sequence can have is sign-flipping
every sample. The link between the two forms is $\alpha = e^{aT_s}$ when $x[n]$
samples $e^{at}$.

\subsection{The sinusoid}
\label{ch02:sinusoid}

The sinusoid is central to the whole book:
\[
  x(t) = A\cos(2\pi f t + \phi) = A\cos(\omega t + \phi),
\]
with \emph{amplitude} $A$, \emph{frequency} $f$ in hertz (so $\omega = 2\pi f$ in
radians/second is the \emph{angular} frequency), and \emph{phase} $\phi$ in
radians. Its discrete cousin is
\[
  x[n] = A\cos(\hat\omega\, n + \phi), \qquad \hat\omega = \omega T_s = 2\pi f / f_s,
\]
where $\hat\omega$ is the \emph{normalized} (or digital) angular frequency, in
radians \emph{per sample}. The sinusoid matters above all others because, as later
chapters will show, it is the one signal shape that passes through any linear
time-invariant system unchanged except in amplitude and phase. That is what makes
``frequency'' the natural language for such systems.

\begin{figure}[h]
\centering
\begin{tikzpicture}
\begin{axis}[title style={font=\small},
  name=ax1, width=0.33\textwidth, height=3.6cm,
  title={$\delta[n]$}, title style={font=\small},
  xmin=-3.5,xmax=3.5,ymin=-0.3,ymax=1.3, axis lines=left,
  xtick={-3,0,3}, ytick={0,1},
  ycomb, mark=*, mark size=1.4pt, color=sigblue, thick]
\addplot coordinates {(-3,0)(-2,0)(-1,0)(0,1)(1,0)(2,0)(3,0)};
\end{axis}
\begin{axis}[title style={font=\small},
  name=ax2, at={(0.33\textwidth,0)}, width=0.33\textwidth, height=3.6cm,
  title={$u[n]$}, title style={font=\small},
  xmin=-3.5,xmax=3.5,ymin=-0.3,ymax=1.3, axis lines=left,
  xtick={-3,0,3}, ytick={0,1},
  ycomb, mark=*, mark size=1.4pt, color=sigblue, thick]
\addplot coordinates {(-3,0)(-2,0)(-1,0)(0,1)(1,1)(2,1)(3,1)};
\end{axis}
\begin{axis}[title style={font=\small},
  name=ax3, at={(0.66\textwidth,0)}, width=0.33\textwidth, height=3.6cm,
  title={$r[n]=n\,u[n]$}, title style={font=\small},
  xmin=-3.5,xmax=3.5,ymin=-0.3,ymax=3.3, axis lines=left,
  xtick={-3,0,3}, ytick={0,1,2,3},
  ycomb, mark=*, mark size=1.4pt, color=sigblue, thick]
\addplot coordinates {(-3,0)(-2,0)(-1,0)(0,0)(1,1)(2,2)(3,3)};
\end{axis}
\end{tikzpicture}

\vspace{0.4em}

\begin{tikzpicture}
\begin{axis}[title style={font=\small},
  width=0.5\textwidth, height=3.8cm,
  title={decaying exponential $\alpha^n u[n]$, $\alpha=0.7$},
  title style={font=\small},
  xmin=-1.5,xmax=8.5,ymin=-0.2,ymax=1.2, axis lines=left,
  ycomb, mark=*, mark size=1.3pt, color=sigblue, thick]
\addplot coordinates {(-1,0)(0,1)(1,0.7)(2,0.49)(3,0.343)(4,0.24)(5,0.168)(6,0.118)(7,0.082)(8,0.058)};
\end{axis}
\begin{axis}[at={(0.5\textwidth,0)}, width=0.5\textwidth, height=3.8cm,
  title={sinusoid $\cos(\hat\omega n)$, $\hat\omega=2\pi/12$},
  title style={font=\small},
  xmin=-0.5,xmax=24.5,ymin=-1.3,ymax=1.3, axis lines=left,
  ycomb, mark=*, mark size=1.1pt, color=sigblue, thick,
  samples at={0,...,24}]
\addplot {cos(deg(2*pi*x/12))};
\end{axis}
\end{tikzpicture}
\caption{The basic discrete-time signals: impulse, step, ramp (top), and the
decaying exponential and a sampled sinusoid (bottom). Every signal we analyze
later is, in some sense, built from these.}
\label{ch02:fig-basic}
\end{figure}

\subsection{The complex exponential -- a sneak preview}

The single most important signal in the book is one we cannot yet fully justify:
\[
  x(t) = e^{\jj \omega t} = \cos(\omega t) + \jj\sin(\omega t),
\]
where $\jj = \sqrt{-1}$ is the imaginary unit (customary in engineering to keep $i$ free for electric current).
Euler's formula $e^{\jj \theta} = \cos\theta + \jj\sin\theta$ shows that $e^{\jj\omega t}$ has unit magnitude
$\abs{e^{\jj\omega t}} = \sqrt{\cos^2(\omega t) + \sin^2(\omega t)} = 1$ at every instant.
Geometrically, $e^{\jj\omega t}$ is a point of unit length rotating around the origin of the complex plane at angular speed $\omega$:
its real part traces a cosine, its imaginary part a sine. It bundles a cosine and a
sine into one rotating arrow, and it is the basic building block of frequency analysis. We mention
it here so the cast of characters is complete; phasors and rotation get the
full geometric treatment in Chapter~4.

\section{Energy and Power}
\label{ch02:energy-power}

How ``big'' is a signal? Two measures dominate, and choosing between them is really
a question about whether your signal is a transient burst or an endless hum.

The instantaneous power of a signal is $\abs{x(t)}^2$ (we use $\abs{\cdot}^2$ so
the definitions work for complex signals too). Integrating that over all time gives
the \textbf{total energy}:
\[
  E = \int_{-\infty}^{\infty} \abs{x(t)}^2 \,\dd t,
  \qquad
  E = \sum_{n=-\infty}^{\infty} \abs{x[n]}^2 .
\]
A signal with finite, nonzero energy is an \emph{energy signal} -- think of a
single clap, a flick of the wrist, anything that starts, happens, and dies away.

But many useful signals never die away: a pure tone, a power-line hum, a constant.
Their energy is infinite, so energy is useless for comparing them. For these we use
average power -- energy per unit time, in the limit:
\[
  P = \lim_{T\to\infty}\frac{1}{2T}\int_{-T}^{T}\abs{x(t)}^2\,\dd t,
  \qquad
  P = \lim_{N\to\infty}\frac{1}{2N+1}\sum_{n=-N}^{N}\abs{x[n]}^2 .
\]
A signal with finite, nonzero average power is a \emph{power signal}. The two
categories are mutually exclusive: an energy signal has $P = 0$ (its finite energy,
spread over infinite time, averages to nothing), and a power signal has
$E = \infty$.

\begin{example}[Power of a sinusoid]
\label{ch02:ex-sine-power}
Take $x(t) = A\cos(\omega t + \phi)$. Using $\cos^2\theta = \tfrac12(1+\cos 2\theta)$,
the $\cos 2\theta$ term averages to zero over any whole number of periods, leaving
\[
  P = \frac{1}{T_0}\int_0^{T_0} A^2\cos^2(\omega t + \phi)\,\dd t
    = \frac{A^2}{2}.
\]
So the average power of a sinusoid is half its amplitude squared, regardless of
frequency or phase. The corresponding root-mean-square (RMS) amplitude is
$A/\sqrt{2}$ -- the number your multimeter reports for an AC voltage. Mains
electricity quoted as ``230~V'' is an RMS figure; its peak is $230\sqrt2 \approx
325$~V.
\end{example}

\begin{example}[Energy of a decaying exponential]
\label{ch02:ex-exp-energy}
$x[n] = \alpha^n u[n]$ with $\abs{\alpha}<1$. Using the standard geometric series
sum $\sum_{n=0}^{\infty} r^n = \frac{1}{1-r}$ (valid for any ratio with $\abs{r}<1$, here with $r = \abs{\alpha}^2$),
we find
\[
  E = \sum_{n=0}^{\infty}\abs{\alpha}^{2n} = \frac{1}{1-\abs{\alpha}^2},
\]
a finite number -- it is an energy signal. For $\alpha = 0.7$,
$E = 1/(1-0.49) \approx 1.96$.
\end{example}

\begin{connectionbox}
Energy is just the squared length (norm) of a signal viewed as a vector:
$E = \norm{x}^2 = \inner{x}{x}$. The next chapter makes this literal -- signals
\emph{are} vectors, energy \emph{is} squared length, and Parseval's theorem (later)
will say that this length is the same whether you measure it in the time domain or
the frequency domain. The link to the probability book is just as direct: average
power is the mean square, and for a zero-mean random signal it equals the variance
$\Var(x)$, the foundation of signal-to-noise ratio.
\end{connectionbox}

\section{Symmetry: Even and Odd}
\label{ch02:even-odd}

A signal is \textbf{even} if it is its own mirror image about $t=0$, and
\textbf{odd} if mirroring flips its sign:
\[
  \text{even:}\quad x(-t) = x(t), \qquad\qquad
  \text{odd:}\quad x(-t) = -x(t).
\]
Cosine is even ($\cos(-\theta)=\cos\theta$); sine is odd ($\sin(-\theta)=-\sin\theta$);
the impulse is even; the ramp is neither. An odd signal must pass through the origin
($x(0) = -x(0) \Rightarrow x(0)=0$).

What makes this useful is a clean algebraic fact: \emph{every} signal splits uniquely into an
even part plus an odd part,
\[
  x_e(t) = \tfrac12\big(x(t) + x(-t)\big), \qquad
  x_o(t) = \tfrac12\big(x(t) - x(-t)\big), \qquad
  x = x_e + x_o.
\]
(Check: $x_e + x_o = x$, and $x_e$ is even while $x_o$ is odd.) The split matters
because the Fourier transform treats even and odd parts very differently --
even parts produce real (cosine) spectra, odd parts produce imaginary (sine)
spectra -- so splitting a signal into even and odd parts corresponds directly to
splitting its spectrum into cosine (real) and sine (imaginary) parts. Symmetry in time therefore produces symmetry in
frequency, and exploiting it is how the real-input FFT runs roughly twice as fast.

\begin{example}[Splitting a one-sided exponential]
\label{ch02:ex-evenodd}
Let $x[n] = \alpha^n u[n]$. Its mirror is $x[-n] = \alpha^{-n}u[-n]$. The even part
is $x_e[n] = \tfrac12(\alpha^{n}u[n] + \alpha^{-n}u[-n])$ -- a two-sided symmetric
decay -- and the odd part is the antisymmetric remainder; at $n=0$, $x_o[0]=0$ and
$x_e[0]=1$, recovering $x[0]=1$.
\end{example}

\section{Periodic and Aperiodic}
\label{ch02:periodic}

A continuous signal is \textbf{periodic} if it repeats after some shift $T_0$,
\[
  x(t + T_0) = x(t) \quad \text{for all } t,
\]
and the smallest such $T_0 > 0$ is the \emph{fundamental period}; its reciprocal
$f_0 = 1/T_0$ is the fundamental frequency. A discrete signal is periodic if
$x[n+N] = x[n]$ for some integer period $N$.

Here the two worlds genuinely diverge, and the difference trips up beginners. A
continuous sinusoid $\cos(\omega t)$ is \emph{always} periodic. A discrete sinusoid
$\cos(\hat\omega n)$ is periodic \emph{only if} $\hat\omega/2\pi$ is a rational
number. The reason: we need $\hat\omega N = 2\pi k$ for integers $N, k$, i.e.
$\hat\omega/2\pi = k/N$ must be rational. If $\hat\omega/2\pi$ is irrational, the
samples never line up to repeat, even though they came from a perfectly periodic
continuous signal.

\begin{example}[A discrete sinusoid that is \emph{not} periodic]
\label{ch02:ex-aperiodic}
$\cos(n)$ -- that is, $\hat\omega = 1$ radian/sample -- has $\hat\omega/2\pi =
1/(2\pi)$, which is irrational. So the sequence $\cos(0), \cos(1), \cos(2), \dots$
never exactly repeats: $\cos(n)$ is an aperiodic discrete signal. By contrast
$\cos(2\pi n/12)$ has $\hat\omega/2\pi = 1/12$, rational, so it repeats every
$N=12$ samples (visible in Figure~\ref{ch02:fig-basic}).
\end{example}

\begin{walkbox}
You sample a 50~Hz tone at $f_s = 200$~Hz. Is the resulting sequence periodic, and
if so, what is its period in samples? \\[0.3em]
\emph{Work it.} $\hat\omega = 2\pi f/f_s = 2\pi (50)/200 = 2\pi/4$, so
$\hat\omega/2\pi = 1/4$: rational, period $N = 4$ samples. Indeed four samples of a
50~Hz tone at 200~Hz span exactly one cycle ($4 \times 5\,\text{ms} = 20\,\text{ms}
= 1/50\,\text{s}$). Now try 60~Hz at 200~Hz: $\hat\omega/2\pi = 60/200 = 3/10$,
period $N=10$.
\end{walkbox}

\begin{trapbox}
A real recorded signal is never truly periodic -- it started at some finite time and
will stop. ``Periodic'' is an idealization we apply to a steady middle stretch.
This matters for the DFT later: it implicitly assumes the finite chunk you fed it
repeats forever, and the mismatch at the seam causes \emph{spectral leakage}, fixed
by windowing. File that away; it is one of the most common surprises in practical
spectrum analysis.
\end{trapbox}

\section{Transformations of the Independent Variable}
\label{ch02:transforms}

Many operations reshape a signal by doing something to its \emph{time axis} rather
than its values. Three are fundamental: shifting, reflecting, and scaling. Getting
the direction of each right is a classic source of sign errors, so we go slowly.

\subsection{Time shift (delay and advance)}

Replacing $t$ by $t - t_0$ \emph{delays} the signal by $t_0$:
\[
  y(t) = x(t - t_0).
\]
It looks backwards but is right: the value $y$ takes at time $t_0$ is the value $x$
took at time $0$, so the whole signal has moved $t_0$ to the right (later). A
positive $t_0$ is a delay; a negative $t_0$ is an advance (the signal moves left,
earlier). In discrete time, $y[n] = x[n - n_0]$ shifts the sequence right by $n_0$
samples. The one-sample delay $x[n-1]$ -- ``yesterday's value'' -- is the core
building block of every digital filter; in hardware it is literally one memory cell.

\begin{intuitionbox}
Why does ``$-t_0$'' move things \emph{right}? Track one feature, say a peak that
$x$ has at $t=0$. In $y(t)=x(t-t_0)$, that peak now occurs where the argument
$t - t_0 = 0$, i.e. at $t = t_0$. The feature relocated from $0$ to $t_0$: rightward
by $t_0$. Subtracting inside the argument adds to the time of every feature.
\end{intuitionbox}

\subsection{Time reversal (reflection)}

$y(t) = x(-t)$ flips the signal left-to-right about $t=0$ -- running the tape
backwards. We already used it to define even/odd. In discrete time, $y[n]=x[-n]$.
Reflection is the operation hiding inside correlation and inside convolution (where
one of the two signals is flipped before sliding), so it is worth being fluent with.

\subsection{Time scaling}

$y(t) = x(at)$ with $a > 1$ \emph{compresses} the signal toward the origin (it plays
faster, finishes sooner); $0 < a < 1$ \emph{stretches} it (plays slower). Again it
can feel backwards: $a=2$ \emph{shrinks} the time axis because a feature at
$x$-time $t_0$ now appears at $y$-time $t_0/2$. A negative $a$ scales \emph{and}
reflects.

Time scaling is where continuous and discrete part ways most sharply. In continuous
time you can scale by \emph{any} real factor. In discrete time, $x[an]$ only makes
sense for integer $an$. Integer compression ($x[2n]$ keeps every other sample --
this is \emph{downsampling}/decimation) and integer expansion (insert zeros between
samples -- \emph{upsampling}) are genuinely different operations from continuous
scaling, with their own subtleties: they can create or destroy frequency content.
Chapter~22 treats them properly.

\begin{example}[Reading off a compound transform]
\label{ch02:ex-compound}
Sketch $y(t) = x(2t - 4)$ given $x(t)$. Mixing up the order of operations is exactly
the trap, so use a fixed recipe: rewrite the argument in the canonical form
$a(t - t_d)$. Here $2t - 4 = 2(t - 2)$, so $a = 2$ and $t_d = 2$. The recipe is
shift-then-scale: take $x(t)$, shift right by $2$ to get $x(t-2)$, then compress the
time axis by $2$ to get $x(2(t-2)) = x(2t-4)$. Verify with one landmark: $y(2)$
should equal $x(0)$, and indeed $2\cdot 2 - 4 = 0$.~\checkmark
\end{example}

\begin{trapbox}
When a transform mixes shift and scale, do not eyeball it. Rewrite the argument in
the canonical form $a(t - t_d)$, read off the scale $a$ and the shift $t_d$, and
apply shift-then-scale. Then verify with a single landmark point (where the
argument is zero). Two minutes of this beats an hour debugging a flipped filter.
\end{trapbox}

\section{The Sifting Property: Building Signals from Impulses}
\label{ch02:sifting}

We close with an idea used throughout the rest of the book. Return to the
discrete impulse. Shift it: $\delta[n-k]$ is a spike sitting at $n=k$. Scale that
spike by a number $c$: $c\,\delta[n-k]$ is a spike of height $c$ at position $k$.
Now add up a whole lineup of such spikes, one for each integer $k$, with the height
of the $k$-th spike chosen to be the actual signal value $x[k]$:
\[
  \boxed{\,x[n] \;=\; \sum_{k=-\infty}^{\infty} x[k]\,\delta[n-k]\,}.
\]
Read it carefully. On the right, for any particular $n$, every term vanishes except
the one with $k = n$ (because $\delta[n-k]=0$ unless $k=n$), and that surviving term
is $x[n]\cdot 1 = x[n]$. So the equation is true: \textbf{any} discrete signal is a sum of shifted, scaled impulses. The
impulse is a \emph{basis} for all signals: a complete set of building blocks, one
per time slot. Figure~\ref{ch02:fig-sifting} shows the decomposition for a short
sequence.

\begin{figure}[h]
\centering
\begin{tikzpicture}
\begin{axis}[title style={font=\small},
  name=sa, width=0.5\textwidth, height=4.0cm,
  title={a short signal $x[n]$}, title style={font=\small},
  xlabel={$n$}, xmin=-1.5,xmax=4.5,ymin=-1.6,ymax=2.6, axis lines=left,
  xtick={-1,0,1,2,3,4}, ytick={-1,0,1,2},
  ycomb, mark=*, mark size=1.6pt, color=sigblue, thick]
\addplot coordinates {(-1,0)(0,2)(1,1)(2,-1)(3,2)(4,0)};
\end{axis}
\begin{axis}[title style={font=\small},
  name=sb, at={(0.5\textwidth,0)}, width=0.5\textwidth, height=4.0cm,
  title={as stacked impulses $x[k]\,\delta[n-k]$},
  title style={font=\small},
  xlabel={$n$}, xmin=-1.5,xmax=4.5,ymin=-1.6,ymax=2.6, axis lines=left,
  xtick={-1,0,1,2,3,4}, ytick={-1,0,1,2}]
\addplot[ycomb, mark=*, mark size=1.6pt, color=specred, thick]
  coordinates {(0,2)};
\addplot[ycomb, mark=*, mark size=1.6pt, color=quotientorange, thick]
  coordinates {(1,1)};
\addplot[ycomb, mark=*, mark size=1.6pt, color=targetgreen, thick]
  coordinates {(2,-1)};
\addplot[ycomb, mark=*, mark size=1.6pt, color=structurepurple, thick]
  coordinates {(3,2)};
\end{axis}
\end{tikzpicture}
\caption{The sifting decomposition. The sequence on the left (blue) is the sum of
four scaled, shifted impulses (right): height $2$ at $n=0$, height $1$ at $n=1$,
height $-1$ at $n=2$, height $2$ at $n=3$. Each impulse carries exactly one sample
value, sitting in its own time slot.}
\label{ch02:fig-sifting}
\end{figure}

\begin{happeningbox}
The sifting sum decomposes a signal by \emph{location in time}: it says ``$x$ is
this much at slot 0, plus this much at slot 1, plus \dots,'' each ``this much''
carried by its own impulse. That sounds almost too obvious to be useful.

The payoff arrives when we study linear systems in Chapters 12 and 13. If you know
how a system responds to a single impulse $\delta[n]$ -- an output sequence we will call
$h[n]$ -- then any input $x[n]$ can be broken into spikes, pushed through the system
spike by spike, and reassembled into the \emph{convolution sum}:
\[
  y[n] = \sum_{k} x[k]\,h[n-k].
\]
A single test spike tells you the output for \emph{every} possible input.
That whole theory starts from the sifting property.
\end{happeningbox}

The continuous version replaces the sum with an integral and the Kronecker delta
with the Dirac delta:
\[
  x(t) = \int_{-\infty}^{\infty} x(\tau)\,\delta(t-\tau)\,\dd\tau,
\]
which is exactly the sifting property of Section~\ref{ch02:impulse} read as ``$x$ is
a continuum of infinitesimal impulses, one at each instant $\tau$, weighted by
$x(\tau)$.''

\begin{whybox}
Why phrase the obvious fact ``$x[n]$ equals $x[n]$'' as an impulse sum at all?
Because it converts the question ``what does a system do to an arbitrary signal?''
into the far smaller question ``what does it do to one impulse?'' The trick is to
write a complicated input as a combination of simple inputs you already understand,
push each simple piece through, and add up the results. We will use this approach
twice more in this book: once with impulses (giving convolution and the
time-domain view) and once with sinusoids/complex exponentials (giving the Fourier
transform and the frequency-domain view). The two choices of building block lead to two
theories, one built on impulses and one built on sinusoids.
\end{whybox}

\begin{goalbox}
For the wrist tracker, the impulse-response idea is how we will reason about the
filters that clean up the IMU stream. The hand motion we care about is slow (a few
hertz); sensor noise and physiological tremor are faster. A smoothing filter is
fully described by its impulse response $h[n]$ -- literally, ``what comes out if you
feed in a single spike.'' Once we have $h[n]$ for our chosen smoother, convolution
tells us what it does to a real gesture, and the Fourier view (later) tells us
\emph{which} frequencies it keeps or kills. The $\delta[n]$ of this
chapter is the test signal we will apply to every filter.
\end{goalbox}

\begin{aestheticbox}
The pieces of this chapter fit together. After separating continuous and
discrete time, introducing the basic elementary signals, and defining energy, power,
symmetry, periodicity, and time-axis shifts, we saw via the sifting property that
one signal in that list -- the unit impulse -- serves as a complete \emph{basis} in
time: every discrete sequence is a weighted sum of shifted impulses. Expressing
arbitrary signals in terms of simple building blocks is the strategy we will use
across the next several chapters.
\end{aestheticbox}

\section{A Tiny Bit of Code}
\label{ch02:code}

Everything above is concrete enough to type. Here is much of the chapter in NumPy:
build the basic sequences, sample a sinusoid, compute energy and power, and verify
the even/odd split.

\begin{lstlisting}[language=Python]
import numpy as np

fs   = 200.0                 # sampling rate, Hz
Ts   = 1.0 / fs
n    = np.arange(0, 200)     # sample indices 0..199 (one second)
t    = n * Ts                # the real time of each sample

# basic signals as sequences
delta = (n == 0).astype(float)          # unit impulse
step  = (n >= 0).astype(float)          # unit step (all ones here)
ramp  = n * step                        # unit ramp
expo  = 0.97 ** n                       # decaying exponential
sine  = np.cos(2*np.pi*5*t)             # 5 Hz sinusoid, sampled

# energy and (average) power of the sinusoid
E = np.sum(np.abs(sine)**2)             # finite-window energy
P = np.mean(np.abs(sine)**2)            # approx A^2/2 = 0.5 for A=1
print(f"energy={E:.2f}  power={P:.3f}  (expect ~0.5)")

# even/odd split of a one-sided exponential
m   = np.arange(-50, 51)
x   = np.where(m >= 0, 0.97**m, 0.0)
xe  = 0.5 * (x + x[::-1])               # even part
xo  = 0.5 * (x - x[::-1])               # odd part
assert np.allclose(xe + xo, x)         # the split reconstructs x
\end{lstlisting}

\begin{labbox}
Run the snippet, then change one thing at a time. (1) Set \texttt{sine =
np.cos(2*np.pi*190*t)} -- a 190~Hz tone sampled at 200~Hz -- and \emph{plot} it with
\texttt{matplotlib}. It will look like a slow 10~Hz wave, not a fast one: your first
encounter with aliasing when sampling too slowly. (2) Change the decay base from
$0.97$ to $1.01$ and watch \texttt{expo} blow up -- crossing $\abs{\alpha}=1$ takes
you from finite-energy decay ($\abs{\alpha}<1$) and constant-power oscillation
($\abs{\alpha}=1$) straight into unbounded growth. (3) Build a delayed impulse
\texttt{np.roll(delta, 5)} and confirm it is $\delta[n-5]$.
\end{labbox}

\begin{historybox}
The elementary signals of this chapter -- the step $u[n]$, the impulse $\delta[n]$, and
discrete sinusoids -- were born directly from the operational realities of early electrical
engineering. In the late 1880s, \textbf{Oliver Heaviside}, working in isolation without academic
backing, sought to calculate the transient surge of voltage when a telegraph switch closed on a
long transmission line. To model the sudden turn-on, Heaviside invented the unit step function
$H(t)$ (today's $u(t)$) and its derivative, the impulsive surge $\dd H/\dd t$.

To solve the differential equations of cables without evaluating laborious integrals, Heaviside
invented his \emph{operational calculus}: he replaced the differential operator
$\dd/\dd t$ with an algebraic symbol $p$, manipulating it as ordinary algebra (even taking square
roots like $\sqrt{p}$ for diffusion lines). The British academic establishment was outraged.
Cambridge mathematicians condemned his work as unprincipled and reckless because he manipulated
divergent operators without proving convergence. When pressed on his lack of formal rigor,
Heaviside famously replied: \emph{``Shall I refuse my dinner because I do not fully
understand the process of digestion?''} Decades later, Norbert Wiener and Laurent Schwartz
formalized Heaviside's operational calculus into modern distribution theory and Laplace transform
analysis.

In continuous time, the impulse $\delta(t)$ was popularized in 1930 by \textbf{Paul Dirac} for
quantum mechanics. Mathematicians balked because a function that is zero everywhere except at the
origin with unit integral cannot exist in classical analysis; Laurent Schwartz finally legitimized
it in 1950 with the theory of distributions (generalized functions defined rigorously by how they
act inside integrals on smooth test functions), earning the Fields Medal. The discrete Kronecker
delta $\delta[n]$ is far tamer, named for \textbf{Leopold Kronecker} (1823--1891). In discrete
time, $\delta[n]$ is simply a legitimate sequence with $\delta[0]=1$ and $\delta[n]=0$ for $n \ne 0$,
free of all continuous analytical paradoxes.
\end{historybox}

\begin{connectionbox}
\textbf{Where We Are and Where We Go Next}
\begin{itemize}
  \item \textbf{What we have so far}: A clean bridge from continuous physical time ($t \in \R$) to
  discrete sample index ($n \in \Z$), together with the sifting identity
  $x[n] = \sum_k x[k]\,\delta[n-k]$ that writes any sequence as a weighted stack of unit impulses.
  \item \textbf{Watch out for}: Discrete time behaves differently from continuous time in two ways.
  First, $\cos(\omega t)$ is periodic for \emph{every} $\omega > 0$, whereas $\cos(\hat{\omega} n)$
  is periodic in $n$ \emph{if and only if} $\hat{\omega}/(2\pi)$ is rational; otherwise the samples
  never repeat. Second, sampling ignores everything that happens between clock ticks, so fast
  oscillations masquerade as slow ones (aliasing) unless we limit the signal's bandwidth before
  sampling.
  \item \textbf{Next}: Chapter~3 views discrete sequences as vectors in $\R^N$ or $\ell^2$, giving
  us lengths, distances, and orthogonal projections. Later, Chapters~12--13 feed the impulse sum
  through linear systems to derive convolution, and Chapter~18 proves the Nyquist--Shannon condition
  for lossless sampling.
\end{itemize}
\end{connectionbox}

\section{Exercises}
\label{ch02:exercises}

\begin{exercise}
A signal is sampled at $f_s = 256$~Hz. (a) What is $T_s$? (b) Which sample index $n$
is closest to $t = 1.0$~s? (c) What real time does sample $n = 400$ correspond to?
\end{exercise}

\begin{exercise}
Classify each as an energy signal, a power signal, or neither, and justify:
(a) $x[n] = (0.5)^n u[n]$; \;(b) $x(t) = 3\cos(20\pi t)$; \;(c) $x[n] = u[n]$ (the
unit step); \;(d) $x(t) = t\,u(t)$ (the ramp).
\end{exercise}

\begin{exercise}
Compute the total energy of $x[n] = (0.8)^n u[n]$. Then compute the energy of the
\emph{delayed} version $x[n-10]$. Does delaying a signal change its energy? Explain
in one sentence why this had to be true.
\end{exercise}

\begin{exercise}
Let $x[n] = \{\,\underline{2}, 3, -1\,\}$ for $n = 0,1,2$ (underline marks $n=0$),
and zero elsewhere. Write $x[n]$ explicitly as a sum of shifted, scaled impulses
$\sum_k x[k]\,\delta[n-k]$.
\end{exercise}

\begin{exercise}
Find the even and odd parts of $x[n] = \{\,\underline{1}, 2, 3\,\}$ for $n=0,1,2$
(zero elsewhere). Verify that $x_e$ is even, $x_o$ is odd, and $x_e + x_o = x$ at
every index.
\end{exercise}

\begin{exercise}
For each discrete sinusoid, state whether it is periodic, and if so give the
fundamental period $N$ in samples: (a) $\cos(2\pi n/8)$; \;(b) $\cos(0.3\pi n)$;
\;(c) $\cos(3 n)$; \;(d) $\cos(2\pi n / 5) + \cos(2\pi n / 3)$ (think about the
least common multiple of the two periods).
\end{exercise}

\begin{exercise}
A continuous sinusoid $\cos(2\pi f t)$ is sampled at $f_s$. Show that the resulting
sequence is periodic in $n$ if and only if $f/f_s$ is rational, and that the period
is then $N = f_s / \gcd(f, f_s)$ when $f, f_s$ are integers. Test it on $f = 60$,
$f_s = 200$.
\end{exercise}

\begin{exercise}
Given an arbitrary $x(t)$, sketch (in words or rough pictures) the result of each
transform, and identify the canonical shift/scale: (a) $x(t+3)$; \;(b) $x(-t)$;
\;(c) $x(3t)$; \;(d) $x(-2t + 6)$. For (d), give the landmark check at the point
where the argument is zero.
\end{exercise}

\begin{exercise}
Using only the relations $\delta[n] = u[n] - u[n-1]$ and $u[n] = \sum_{k\le n}
\delta[k]$, prove that the ramp satisfies $r[n] - r[n-1] = u[n-1]$. (Hint: write the
ramp as a running sum of the step and difference it.)
\end{exercise}

\begin{exercise}
\emph{Toward the wrist tracker.} Suppose your IMU streams acceleration at
$f_s = 200$~Hz, and a typical wrist ``flick'' gesture lasts about 150~ms. (a) How
many samples does one flick span? (b) If you instead sampled at only $f_s = 20$~Hz,
how many samples would the flick span, and why might that be too few to recognize
the gesture reliably? (c) Argue qualitatively, using the energy idea, why a brief
flick is an ``energy signal'' while the steady hum of holding your arm still
resembles a ``power signal'' plus a constant offset (gravity).
\end{exercise}

\begin{exercise}
A digital delay element outputs $y[n] = x[n-1]$. Its response to a single impulse $\delta[n]$
is $h[n] = \delta[n-1]$. Using the sifting decomposition $x[n] = \sum_k x[k]\,\delta[n-k]$,
confirm that delaying each spike by one sample reassembles into $\sum_k x[k]\,\delta[n-1-k] = x[n-1]$.
This previews why representing signals as sums of impulses will let us calculate the effect
of any linear filter once we reach Chapter 13.
\end{exercise}
